\documentclass[a4paper]{article}
\usepackage[utf8]{inputenc}
\usepackage[T5]{fontenc}
\usepackage[vietnamese]{babel}
\usepackage[fontsize=13pt]{scrextend}
\usepackage{mathptmx,amsmath}
\usepackage{geometry,setspace,indentfirst}
\usepackage{graphicx,xcolor,booktabs,longtable,tabularx,array}
\usepackage{enumitem,listings,caption,float,fancyhdr,hyperref,microtype}
\usepackage{tikz}
\usetikzlibrary{arrows.meta,positioning,fit,shapes.geometric,calc}

\geometry{left=3cm,right=2cm,top=2cm,bottom=2cm}
\setstretch{1.3}
\setlength{\parindent}{1cm}
\setlength{\parskip}{3pt}
\setlength{\emergencystretch}{3em}
\setlist{nosep}
\renewcommand{\arraystretch}{1.25}

\definecolor{bluec}{RGB}{24,91,166}
\definecolor{greenc}{RGB}{28,132,84}
\definecolor{orangec}{RGB}{210,116,24}
\definecolor{redc}{RGB}{183,54,54}
\definecolor{violetc}{RGB}{105,75,170}
\definecolor{grayc}{RGB}{92,100,112}
\definecolor{codebg}{RGB}{247,248,250}
\hypersetup{colorlinks=true,linkcolor=bluec,urlcolor=bluec,
  pdftitle={Báo cáo DashboardManagerWifiMesh},pdfauthor={Nhóm phát triển MSMS}}

\lstdefinestyle{textstyle}{
  backgroundcolor=\color{codebg},frame=single,rulecolor=\color{gray!35},
  basicstyle=\ttfamily\small,breaklines=true,columns=fullflexible,
  showstringspaces=false
}
\pagestyle{fancy}
\fancyhf{}
\fancyhead[L]{Báo cáo DashboardManagerWifiMesh}
\fancyhead[R]{React / Node.js}
\fancyfoot[C]{\thepage}
\setlength{\headheight}{16pt}

\newcommand{\codeword}[1]{\texttt{\detokenize{#1}}}
\newcommand{\pathcode}[1]{\nolinkurl{#1}}
\tikzset{
  block/.style={draw=bluec,rounded corners=3pt,thick,align=center,
    minimum height=9mm,fill=blue!5,font=\small},
  service/.style={draw=violetc,rounded corners=3pt,thick,align=center,
    minimum height=9mm,fill=violet!7,font=\small},
  store/.style={draw=orangec,thick,align=center,fill=orange!9,
    cylinder,shape border rotate=90,aspect=.25,font=\small},
  decision/.style={draw=orangec,thick,diamond,aspect=2,align=center,
    fill=orange!8,font=\small},
  flow/.style={-{Latex[length=2.4mm]},thick,draw=bluec},
  dataflow/.style={-{Latex[length=2.4mm]},thick,draw=greenc},
  warnflow/.style={-{Latex[length=2.4mm]},thick,dashed,draw=redc},
  note/.style={draw=grayc,rounded corners=2pt,fill=gray!8,align=left,font=\scriptsize}
}

\begin{document}
\pagenumbering{Alph}
\begin{titlepage}
\centering
\vspace*{1.2cm}
{\large BÁO CÁO PHÂN TÍCH MÃ NGUỒN\par}
\vspace{1.5cm}
{\Huge\bfseries DashboardManagerWifiMesh\par}
\vspace{.8cm}
{\Large Dashboard giám sát, WebSocket server và lưu trữ dữ liệu MSMS\par}
\vspace{2cm}
\begin{tabular}{rl}
Frontend: & React 18, Material UI, ApexCharts\\
Backend: & Node.js, Express, ws\\
Lưu trữ: & SQLite hoặc MongoDB\\
Giao tiếp: & HTTP REST, WebSocket, mDNS\\
Phiên bản tài liệu: & 1.0\\
\end{tabular}
\vfill
{\large Tài liệu được xây dựng từ trạng thái mã nguồn hiện tại\par}
\vspace{.3cm}
{\large Tháng 07 năm 2026\par}
\end{titlepage}

\pagenumbering{roman}
\tableofcontents
\clearpage
\listoffigures
\clearpage
\listoftables
\clearpage
\pagenumbering{arabic}

\section{Giới thiệu}

\codeword{DashboardManagerWifiMesh} là ứng dụng giám sát gồm frontend React và
backend Node.js chạy trong cùng repository. Backend nhận bản tin từ MeshGateWay,
phân tách payload compact của mạng mesh, lưu từng giá trị đo vào cơ sở dữ liệu,
rồi phát bản tin cho các trình duyệt đang mở dashboard. Frontend kết hợp snapshot
HTTP với stream WebSocket để dựng trạng thái gateway, node, cảm biến, lịch sử và
hiệu năng server.

\subsection{Mục tiêu tài liệu}

\begin{itemize}
 \item giải thích đầy đủ đường đi của dữ liệu từ gateway tới từng màn hình;
 \item mô tả parser hai lớp \codeword{uart_rx} và schema compact của firmware;
 \item phân biệt dữ liệu realtime trong RAM với dữ liệu lịch sử trong database;
 \item làm rõ cách chọn local/server/custom WebSocket;
 \item chỉ ra chức năng đã hoạt động, chức năng giao diện mẫu và các điểm không
 tương thích cần sửa.
\end{itemize}

\section{Kiến trúc tổng thể}

\begin{figure}[H]
\centering
\resizebox{.96\textwidth}{!}{
\begin{tikzpicture}[node distance=14mm and 19mm]
 \node[block,minimum width=3.2cm,fill=green!8] (gw)
   {\textbf{MeshGateWay}\\\codeword{uart_rx}\\\codeword{gateway_status}};
 \node[service,right=of gw,minimum width=4.1cm,minimum height=1.7cm] (server)
   {\textbf{Node.js server :9090}\\HTTP + WebSocket\\parser + broadcast};
 \node[store,below=18mm of server,minimum width=3.2cm] (db)
   {SQLite mặc định\\hoặc MongoDB};
 \node[block,right=of server,minimum width=4.1cm,minimum height=1.7cm] (react)
   {\textbf{React :3000}\\MeshRealtimeProvider\\routes + layouts};
 \node[block,right=of react,minimum width=3.0cm,fill=violet!8] (user)
   {\textbf{Trình duyệt}\\Dashboard\\History / Settings};

 \draw[dataflow] (gw) -- node[above,align=center,font=\scriptsize]
   {WebSocket\\local WS / server WSS} (server);
 \draw[dataflow] (server) -- node[right,font=\scriptsize]{persist/query} (db);
 \draw[dataflow] (server) -- node[above,align=center,font=\scriptsize]
   {WS realtime\\REST history} (react);
 \draw[flow] (react) -- (user);
\end{tikzpicture}}
\caption{Kiến trúc triển khai của DashboardManagerWifiMesh}
\label{fig:architecture}
\end{figure}

Trong chế độ development, \codeword{npm run dev} dùng \codeword{concurrently}
để chạy React và WebSocket server. Hai tiến trình không phải một server duy nhất:
React phục vụ trang tại cổng 3000; Node.js lắng nghe HTTP và WebSocket tại cổng
9090.

\section{Tổ chức repository và công nghệ}

\begin{longtable}{p{5.3cm}p{8.7cm}}
\caption{Các thư mục và tệp quan trọng}\\
\toprule
\textbf{Đường dẫn} & \textbf{Trách nhiệm}\\
\midrule\endfirsthead
\toprule
\textbf{Đường dẫn} & \textbf{Trách nhiệm}\\
\midrule\endhead
\pathcode{public/index.html} & Vỏ HTML và điểm mount \codeword{root}.\\
\pathcode{src/index.js} & Bọc Router, UI context, realtime context và App.\\
\pathcode{src/App.js} & Theme, sidenav và runtime routing.\\
\pathcode{src/routes.js} & Bảng route đồng thời là nguồn tạo menu.\\
\pathcode{src/context/meshRealtime.js} & Kết nối WS dùng chung, parser, registry và state realtime.\\
\pathcode{src/utils/wsConfig.js} & Chọn URL local/server/custom và lưu localStorage.\\
\pathcode{src/utils/meshUdpJsonSchema.js} & Chuẩn hóa payload compact, sensor và đơn vị.\\
\pathcode{src/layouts/} & Các trang Dashboard, Gateway, Nodes, Sensors, History, Alerts, OTA và Settings.\\
\pathcode{websocket-server/server.js} & HTTP REST, WS, mDNS, heartbeat, broadcast và server metrics.\\
\pathcode{websocket-server/sqlite-store.js} & SQLite WAL, ghi dữ liệu, history và CSV.\\
\pathcode{websocket-server/mongodb-store.js} & Adapter MongoDB cũ, chưa đồng nhất interface mới.\\
\pathcode{execute.bat / stop.bat} & Khởi chạy và dừng cây tiến trình trên Windows.\\
\bottomrule
\end{longtable}

\begin{table}[H]
\centering
\caption{Thư viện chính}
\begin{tabularx}{\textwidth}{p{4cm}X}
\toprule
\textbf{Thư viện} & \textbf{Vai trò}\\
\midrule
React 18, React Router 5 & Cây component, state và điều hướng SPA.\\
Material UI, Emotion & Component UI, theme, CSS-in-JS và RTL.\\
ApexCharts & Biểu đồ realtime và history.\\
Express, CORS & REST API và middleware HTTP.\\
ws & WebSocket server, client registry, ping/pong.\\
sqlite + sqlite3 & CSDL file local, chế độ WAL.\\
mongoose & Schema và truy vấn MongoDB.\\
systeminformation & CPU, RAM, nhiệt độ và uptime máy chạy backend.\\
bonjour-service & Quảng bá \codeword{systemmsems.local} bằng mDNS.\\
Swagger & Tài liệu REST tại \codeword{/api-docs}.\\
\bottomrule
\end{tabularx}
\end{table}

\section{Khởi động và định tuyến frontend}

\begin{figure}[H]
\centering
\resizebox{.9\textwidth}{!}{
\begin{tikzpicture}[node distance=7mm]
 \node[block,minimum width=8.5cm] (html) {Trình duyệt tải \codeword{public/index.html}};
 \node[block,below=of html,minimum width=8.5cm] (index) {\codeword{src/index.js}: \codeword{createRoot()}};
 \node[block,below=of index,minimum width=8.5cm] (providers)
   {BrowserRouter $\rightarrow$ VisionUIControllerProvider $\rightarrow$ MeshRealtimeProvider};
 \node[block,below=of providers,minimum width=8.5cm] (app)
   {\codeword{App.js}: theme, sidenav, route switch};
 \node[block,below=of app,minimum width=8.5cm] (routes)
   {\codeword{routes.js}: chọn layout theo URL};
 \node[block,below=of routes,minimum width=8.5cm,fill=green!8] (page)
   {Layout đọc dữ liệu chung bằng \codeword{useMeshRealtime()}};
 \foreach \a/\b in {html/index,index/providers,providers/app,app/routes,routes/page}
   \draw[flow] (\a)--(\b);
\end{tikzpicture}}
\caption{Trình tự mount và render frontend}
\end{figure}

\subsection{Các route hiện tại}

\begin{table}[H]
\centering
\caption{Bản đồ màn hình}
\begin{tabularx}{\textwidth}{p{4cm}p{4.2cm}X}
\toprule
\textbf{Route} & \textbf{Layout} & \textbf{Mục đích}\\
\midrule
\codeword{/dashboard} & Dashboard & Tổng quan root, node, throughput và cảm biến mới nhất.\\
\codeword{/gateway} & Gateway & CPU/RAM/RSSI/pin/uptime của MeshGateWay.\\
\codeword{/nodes} & Nodes & Danh sách registry node và trạng thái online.\\
\codeword{/nodes/:nodeId} & NodeDetail & Chi tiết một node; route ẩn khỏi menu.\\
\codeword{/sensors} & Sensors & Các giá trị sensor realtime.\\
\codeword{/history} & History & Chọn MAC/sensor/field, biểu đồ và export CSV.\\
\codeword{/alerts} & Alerts & Hiển thị lỗi runtime và điều kiện cảnh báo.\\
\codeword{/ota/overview} & OtaOverview & Tổng quan OTA.\\
\codeword{/ota/firmware} & FirmwareManagement & Giao diện quản lý firmware.\\
\codeword{/ota/fota} & UpdateNodes & Giao diện FOTA.\\
\codeword{/settings/system} & SystemSettings & Chọn nguồn WS và các cấu hình UI.\\
\codeword{/about} & About & Thông tin hệ thống.\\
\bottomrule
\end{tabularx}
\end{table}

\section{Lựa chọn nguồn WebSocket}

\subsection{Thứ tự xác định URL}

\begin{figure}[H]
\centering
\resizebox{.88\textwidth}{!}{
\begin{tikzpicture}[node distance=8mm]
 \node[decision] (stored) {localStorage có\\URL hợp lệ?};
 \node[block,right=23mm of stored,minimum width=4.5cm,fill=green!8] (useStored)
   {Dùng \codeword{msms_websocket_url}};
 \node[decision,below=16mm of stored] (env) {\codeword{REACT_APP_WS_URL}\\được đặt?};
 \node[block,right=23mm of env,minimum width=4.5cm] (useEnv) {Dùng URL build environment};
 \node[block,below=16mm of env,minimum width=4.5cm,fill=orange!8] (local)
   {Mặc định\\\codeword{ws://<browser-host>:9090/ws}};
 \draw[dataflow] (stored)--node[above,font=\scriptsize]{có}(useStored);
 \draw[flow] (stored)--node[left,font=\scriptsize]{không}(env);
 \draw[dataflow] (env)--node[above,font=\scriptsize]{có}(useEnv);
 \draw[flow] (env)--node[left,font=\scriptsize]{không}(local);
\end{tikzpicture}}
\caption{Thứ tự chọn URL trong wsConfig.js}
\end{figure}

Ba lựa chọn trên giao diện Settings:
\begin{itemize}
 \item \textbf{Local server:} \codeword{ws://<window.location.hostname>:9090/ws}.
 Khi mở frontend bằng \codeword{localhost:3000}, URL là
 \codeword{ws://localhost:9090/ws}.
 \item \textbf{Internet server:}
 \codeword{wss://systemmsems.msems.click/ws}.
 \item \textbf{Custom URL:} người dùng nhập URL bắt đầu bằng \codeword{ws://}
 hoặc \codeword{wss://}.
\end{itemize}

\codeword{saveWebSocketSourceSettings()} kiểm tra URL, ghi source và URL vào
\codeword{localStorage}, sau đó phát custom event
\codeword{msms:websocket-config-changed}. Provider nhận event, đổi dependency
\codeword{url}, đóng socket cũ và mở socket mới. Khác với URL trên firmware
gateway, cấu hình frontend này bền qua refresh trình duyệt nhưng chỉ thuộc profile
trình duyệt hiện tại.

\begin{figure}[H]
\centering
\resizebox{.94\textwidth}{!}{
\begin{tikzpicture}[node distance=10mm and 13mm]
 \node[block,minimum width=3.1cm] (form) {Settings\\Apply};
 \node[block,right=of form,minimum width=3.4cm] (storage) {Validate\\ghi localStorage};
 \node[block,right=of storage,minimum width=3.4cm] (event) {Custom event\\URL changed};
 \node[service,right=of event,minimum width=3.8cm] (provider) {MeshRealtimeProvider\\cleanup socket cũ};
 \node[block,below=13mm of provider,minimum width=3.8cm,fill=green!8] (newws)
   {Mở WebSocket mới\\reconnect mỗi 3 s};
 \draw[flow] (form)--(storage);
 \draw[flow] (storage)--(event);
 \draw[flow] (event)--(provider);
 \draw[dataflow] (provider)--(newws);
\end{tikzpicture}}
\caption{Áp dụng nguồn dữ liệu mới không cần sửa cứng mã}
\end{figure}

\section{WebSocket và HTTP backend}

\subsection{Khởi tạo server}

Backend dùng một \codeword{http.Server} chung cho Express và
\codeword{WebSocketServer}. Mặc định bind \codeword{0.0.0.0:9090}. mDNS quảng bá
\codeword{systemmsems.local}, service HTTP với TXT
\codeword{service=websocket,path=/ws}. WebSocket server không giới hạn path trong
constructor, vì vậy kết nối tới root hoặc \codeword{/ws} đều đi vào cùng server.

\begin{figure}[H]
\centering
\resizebox{.9\textwidth}{!}{
\begin{tikzpicture}[node distance=8mm]
 \node[block,minimum width=8.4cm] (env) {Đọc HOST, PORT, DB\_TYPE, MONGO\_URI, MDNS\_HOSTNAME};
 \node[block,below=of env,minimum width=8.4cm] (store) {Chọn SQLite store hoặc Mongo store};
 \node[block,below=of store,minimum width=8.4cm] (http) {Tạo Express + REST + Swagger + HTTP server};
 \node[block,below=of http,minimum width=8.4cm] (ws) {Gắn WebSocketServer lên HTTP server};
 \node[block,below=of ws,minimum width=8.4cm] (listen) {Listen cổng 9090 và quảng bá mDNS};
 \node[block,below=of listen,minimum width=8.4cm] (connect) {Kết nối database; log lỗi nhưng server vẫn tiếp tục};
 \foreach \a/\b in {env/store,store/http,http/ws,ws/listen,listen/connect}\draw[flow](\a)--(\b);
\end{tikzpicture}}
\caption{Khởi động backend Node.js}
\end{figure}

\subsection{Xử lý một kết nối}

\begin{figure}[H]
\centering
\resizebox{.94\textwidth}{!}{
\begin{tikzpicture}[node distance=8mm and 10mm]
 \node[block,minimum width=3.1cm] (conn) {Client kết nối};
 \node[block,right=of conn,minimum width=3.2cm] (state) {\codeword{isAlive=true}\\\codeword{uartCarry=""}};
 \node[block,right=of state,minimum width=3.2cm] (listeners) {Gắn message, pong, close};
 \node[block,right=of listeners,minimum width=3.3cm,fill=green!8] (welcome)
   {Gửi \codeword{welcome}\\serverNet};
 \node[block,below=15mm of state,minimum width=3.2cm] (heartbeat) {Mỗi 30 s\\ping/pong};
 \node[block,right=of heartbeat,minimum width=3.2cm,fill=red!7] (zombie)
   {Không pong\\terminate};
 \draw[flow](conn)--(state);
 \draw[flow](state)--(listeners);
 \draw[dataflow](listeners)--(welcome);
 \draw[flow](state)--(heartbeat);
 \draw[warnflow](heartbeat)--(zombie);
\end{tikzpicture}}
\caption{Vòng đời kết nối WebSocket phía server}
\end{figure}

\section{Pipeline nhận và phân tách bản tin}

\begin{figure}[H]
\centering
\resizebox{\textwidth}{!}{
\begin{tikzpicture}[node distance=7mm and 7mm]
 \node[block,minimum width=2.3cm] (raw) {WS frame};
 \node[decision,right=of raw] (binary) {Binary?};
 \node[block,above=12mm of binary,minimum width=2.5cm] (bcastbin) {Broadcast nguyên bytes};
 \node[decision,right=of binary] (json) {JSON hợp lệ?};
 \node[block,above=12mm of json,minimum width=2.7cm] (text) {Bọc \codeword{type=text}\\rồi broadcast};
 \node[decision,right=of json] (type) {Loại payload?};
 \node[block,below left=15mm and 6mm of type,minimum width=2.7cm] (uart)
   {\codeword{uart_rx}\\parse payload/hex};
 \node[block,below=15mm of type,minimum width=2.7cm] (mesh)
   {Mesh JSON\\trực tiếp};
 \node[block,below right=15mm and 6mm of type,minimum width=2.8cm] (gateway)
   {\codeword{gateway_status}};
 \node[store,below=16mm of mesh,minimum width=2.8cm] (db) {Persist rows};
 \node[block,right=18mm of gateway,minimum width=2.8cm,fill=green!8] (broadcast)
   {Broadcast JSON gốc\\trừ client gửi};

 \draw[flow](raw)--(binary);
 \draw[dataflow](binary)--node[right,font=\scriptsize]{có}(bcastbin);
 \draw[flow](binary)--node[above,font=\scriptsize]{không}(json);
 \draw[warnflow](json)--node[right,font=\scriptsize]{không}(text);
 \draw[flow](json)--node[above,font=\scriptsize]{có}(type);
 \draw[dataflow](type)--(uart);
 \draw[dataflow](type)--(mesh);
 \draw[dataflow](type)--(gateway);
 \draw[dataflow](uart)--(db);
 \draw[dataflow](mesh)--(db);
 \draw[dataflow](gateway)--(db);
 \draw[dataflow](type)--(broadcast);
\end{tikzpicture}}
\caption{Phân nhánh xử lý message tại backend}
\label{fig:serverpipeline}
\end{figure}

Lỗi parser hoặc database được bắt riêng. Một JSON hợp lệ vẫn được broadcast với
cấu trúc ban đầu ngay cả khi persist thất bại. Đây là quyết định tốt cho realtime:
UI không bị biến bản tin hợp lệ thành text chỉ vì database gặp lỗi.

\subsection{Ba dạng mesh input được chấp nhận}

\begin{enumerate}
 \item JSON mesh trực tiếp chứa \codeword{v}, \codeword{n}, \codeword{M} hoặc
 \codeword{i}, và mảng \codeword{p}.
 \item \codeword{uart_rx.payload} là object hoặc chuỗi JSON.
 \item \codeword{uart_rx.hex} là UTF-8 mã hex; server ghép phần dòng chưa hoàn
 chỉnh trong \codeword{state.uartCarry}, giới hạn phần carry 4096 ký tự.
\end{enumerate}

\begin{figure}[H]
\centering
\resizebox{.93\textwidth}{!}{
\begin{tikzpicture}[node distance=8mm]
 \node[block,minimum width=7.6cm] (outer) {Parse JSON ngoài: \codeword{type=uart_rx}};
 \node[decision,below=of outer] (payload) {\codeword{payload} có sẵn?};
 \node[block,below left=14mm and 22mm of payload,minimum width=4cm] (pobj)
   {Object: kiểm tra schema\\String: parse JSON hoặc từng dòng};
 \node[decision,below right=14mm and 22mm of payload] (hex) {Có \codeword{hex}?};
 \node[block,below=13mm of hex,minimum width=4cm] (decode)
   {Decode UTF-8\\carry + split CR/LF};
 \node[block,below=17mm of payload,minimum width=4cm,fill=green!8] (packets)
   {Danh sách mesh packet hợp lệ};
 \draw[dataflow](outer)--(payload);
 \draw[dataflow](payload)--node[above left,font=\scriptsize]{có}(pobj);
 \draw[flow](payload)--node[above right,font=\scriptsize]{không}(hex);
 \draw[dataflow](hex)--node[right,font=\scriptsize]{có}(decode);
 \draw[dataflow](pobj)--(packets);
 \draw[dataflow](decode)--(packets);
\end{tikzpicture}}
\caption{Parser hai lớp của uart\_rx}
\end{figure}

\section{Schema dữ liệu mesh}

\begin{lstlisting}[style=textstyle,caption={Payload compact tiêu biểu}]
{
  "v": 1,
  "seq": 88,
  "n": 2,
  "M": "B4:BF:E9:34:0C:CC",
  "t": "2012-01-01T00:02:02",
  "ver": "0.0.1",
  "err": ["0x51"],
  "p": [[1, 6, 27.594, 42.450]]
}
\end{lstlisting}

\begin{table}[H]
\centering
\caption{Ý nghĩa header mesh}
\begin{tabularx}{\textwidth}{p{2.2cm}p{3cm}X}
\toprule
\textbf{Trường} & \textbf{Kiểu} & \textbf{Ý nghĩa}\\
\midrule
\codeword{v} & number & Phiên bản schema. Parser frontend khai báo hằng số 2 nhưng không bắt buộc bằng đúng 2.\\
\codeword{seq} & number & Số thứ tự dùng tính packet loss tại backend.\\
\codeword{n} & number & Tầng mesh.\\
\codeword{M} / \codeword{i} & string & Định danh node; mã hiện tại ưu tiên \codeword{M}.\\
\codeword{t} & ISO string & RTC của node, dùng ước lượng latency.\\
\codeword{ver} & string & Phiên bản firmware.\\
\codeword{err} & string array & Mã lỗi dạng \codeword{0xNN}.\\
\codeword{p} & array & Các hàng \codeword{[port,type,v0,...]}.\\
\bottomrule
\end{tabularx}
\end{table}

\subsection{Ánh xạ loại cảm biến}

\begin{table}[H]
\centering
\caption{Sensor type và thứ tự giá trị}
\begin{tabularx}{\textwidth}{p{1.5cm}p{3.4cm}X}
\toprule
\textbf{ID} & \textbf{Sensor} & \textbf{Các trường theo thứ tự}\\
\midrule
0 & BME280 & \codeword{temp_C}, \codeword{pressure_hPa}, \codeword{humidity_RH}\\
1 & MH-Z14A & \codeword{co2_ppm}, \codeword{temp_C}\\
2 & PMS7003 & \codeword{pm1_0_ugm3}, \codeword{pm2_5_ugm3}, \codeword{pm10_ugm3}\\
3 & DHT22 & \codeword{temp_C}, \codeword{humidity_RH}\\
4 & AHT10 & \codeword{temp_C}, \codeword{humidity_RH}\\
5 & DHT11 & \codeword{temp_C}, \codeword{humidity_RH}\\
6 & HTU21D & \codeword{temp_C}, \codeword{humidity_RH}\\
\bottomrule
\end{tabularx}
\end{table}

Frontend giới hạn tối đa năm giá trị trên một row. \codeword{wirePort} hợp lệ từ
1 đến 255 và \codeword{portIndex = wirePort - 1}. Mã
\codeword{isSensorTypeWireValue()} hiện kiểm tra tối đa
\codeword{SENSOR_AHT10} (ID 4), trong khi registry định nghĩa cả DHT11 và HTU21D;
đây là điểm không đồng nhất cần sửa nếu hàm kiểm tra này được dùng để reject type.

\section{Lưu trữ SQLite}

\subsection{Mô hình dữ liệu dạng long table}

Mỗi giá trị sensor trở thành một dòng. Ví dụ một row BME280 có ba giá trị sẽ tạo
ba record cùng MAC, thời gian, port và sensorName nhưng khác \codeword{field}.

\begin{figure}[H]
\centering
\resizebox{.92\textwidth}{!}{
\begin{tikzpicture}[node distance=9mm]
 \node[block,minimum width=7.5cm] (packet)
   {Packet \codeword{p=[[1,0,27.1,1007.2,55.0]]}};
 \node[block,below left=13mm and 25mm of packet,minimum width=3.5cm] (r1)
   {field=temp\_C\\value=27.1};
 \node[block,below=13mm of packet,minimum width=3.5cm] (r2)
   {field=pressure\_hPa\\value=1007.2};
 \node[block,below right=13mm and 25mm of packet,minimum width=3.5cm] (r3)
   {field=humidity\_RH\\value=55.0};
 \node[store,below=18mm of r2,minimum width=4cm] (data) {Bảng Data};
 \draw[dataflow](packet)--(r1);
 \draw[dataflow](packet)--(r2);
 \draw[dataflow](packet)--(r3);
 \draw[dataflow](r1)--(data);
 \draw[dataflow](r2)--(data);
 \draw[dataflow](r3)--(data);
\end{tikzpicture}}
\caption{Chuẩn hóa một packet thành nhiều record SQLite}
\end{figure}

\begin{longtable}{p{3.2cm}p{2.8cm}p{8cm}}
\caption{Cột bảng Data của SQLite}\\
\toprule
\textbf{Cột} & \textbf{Kiểu} & \textbf{Nội dung}\\
\midrule\endfirsthead
\toprule
\textbf{Cột} & \textbf{Kiểu} & \textbf{Nội dung}\\
\midrule\endhead
\codeword{id} & INTEGER & Khóa chính tự tăng.\\
\codeword{mac} & TEXT & \codeword{M} hoặc \codeword{i}.\\
\codeword{deviceTime} & TEXT & Chuỗi thời gian thiết bị.\\
\codeword{serverTime} & TEXT & ISO time tại lúc backend nhận.\\
\codeword{meshLevel} & INTEGER & Tầng mesh.\\
\codeword{schemaVersion} & INTEGER & Trường \codeword{v}.\\
\codeword{packetLoss} & REAL & Tỷ lệ mất gói nếu có.\\
\codeword{port} & INTEGER & Port trên wire.\\
\codeword{sensorType} & INTEGER & ID sensor.\\
\codeword{sensorName} & TEXT & Tên đã ánh xạ.\\
\codeword{field} & TEXT & Tên đại lượng.\\
\codeword{value} & REAL & Giá trị số.\\
\codeword{nodeTimeValid} & INTEGER & Hiện được ghi 0 trong adapter SQLite.\\
\codeword{sourceIp} & TEXT & IP socket gửi tới backend.\\
\bottomrule
\end{longtable}

SQLite bật \codeword{PRAGMA journal_mode=WAL}, tạo index theo
\codeword{serverTime DESC} và index ghép
\codeword{(mac,field,serverTime)}. Chế độ \codeword{execute.bat} đặt
\codeword{DB_TYPE=sqlite}; database mặc định là
\pathcode{websocket-server/mesh-data.sqlite}, kèm các file WAL/SHM khi đang chạy.

\subsection{Lưu gateway\_status}

Adapter SQLite chuyển năm metrics thành record có
\codeword{sensorName=Gateway}: CPU load, RAM used, RSSI, battery voltage và
uptime. MAC gateway được lấy lần lượt từ \codeword{mac},
\codeword{sta_mac}, \codeword{sta_ip}, source IP hoặc chuỗi
\codeword{gateway}.

\section{REST API và dữ liệu lịch sử}

\begin{table}[H]
\centering
\caption{REST endpoint}
\begin{tabularx}{\textwidth}{p{5.1cm}X}
\toprule
\textbf{Endpoint} & \textbf{Kết quả}\\
\midrule
\codeword{GET /api/nodes} & Snapshot node theo record mới nhất.\\
\codeword{GET /api/history/meta?limit=} & Danh sách MAC, lastSeen và sensor names.\\
\codeword{GET /api/history/series} & Chuỗi tối đa 5000 điểm theo MAC, sensor, field và thời gian.\\
\codeword{GET /api/history/export} & Stream CSV toàn bộ history của một MAC theo khoảng thời gian.\\
\codeword{GET /api-docs} & Swagger UI.\\
\bottomrule
\end{tabularx}
\end{table}

\begin{figure}[H]
\centering
\resizebox{.92\textwidth}{!}{
\begin{tikzpicture}[node distance=8mm and 11mm]
 \node[block,minimum width=3.4cm] (history) {Trang History\\chọn MAC/field/time};
 \node[service,right=of history,minimum width=3.7cm] (api) {Express REST\\validate query};
 \node[store,right=of api,minimum width=3.5cm] (db) {SQLite Data};
 \node[block,below=14mm of api,minimum width=3.7cm] (response)
   {JSON timeseries\\hoặc CSV stream};
 \draw[flow](history)--node[above,font=\scriptsize]{HTTP GET}(api);
 \draw[dataflow](api)--(db);
 \draw[dataflow](db)--(response);
 \draw[dataflow](response)-|(history);
\end{tikzpicture}}
\caption{Luồng truy vấn lịch sử}
\end{figure}

Export SQLite đọc theo chunk 5000 dòng và ghi trực tiếp vào HTTP response, tránh
nạp toàn bộ file CSV vào RAM. Tuy nhiên nội dung CSV được ghép bằng template
string mà chưa escape dấu phẩy, dấu ngoặc kép hay xuống dòng; tên sensor hiện tại
an toàn nhưng nên dùng hàm CSV escaping chuẩn.

\section{Realtime state ở frontend}

\subsection{MeshRealtimeProvider}

\begin{figure}[H]
\centering
\resizebox{.96\textwidth}{!}{
\begin{tikzpicture}[node distance=12mm and 16mm]
 \node[service,minimum width=4cm,minimum height=2.1cm] (provider)
   {\textbf{MeshRealtimeProvider}\\một kết nối WS dùng chung\\một registry node};
 \node[block,above left=of provider,minimum width=3.4cm] (http) {HTTP snapshot\\\codeword{/api/nodes}};
 \node[block,above right=of provider,minimum width=3.4cm] (ws) {WebSocket\\realtime};
 \node[block,below left=of provider,minimum width=3.4cm] (state1)
   {gatewayStatus\\serverMetrics\\latestSensors};
 \node[block,below=of provider,minimum width=3.4cm] (state2)
   {registry / nodes\\online/offline};
 \node[block,below right=of provider,minimum width=3.4cm] (state3)
   {throughput\\latency / logs};
 \node[block,right=24mm of provider,minimum width=3.4cm,fill=green!8] (layouts)
   {Dashboard\\Gateway\\Nodes\\Sensors};
 \draw[dataflow](http)--(provider);
 \draw[dataflow](ws)--(provider);
 \draw[dataflow](provider)--(state1);
 \draw[dataflow](provider)--(state2);
 \draw[dataflow](provider)--(state3);
 \draw[flow](provider)--(layouts);
\end{tikzpicture}}
\caption{State realtime được chia sẻ cho toàn bộ layout}
\end{figure}

Provider fetch snapshot một lần khi URL thay đổi, sau đó mở WebSocket. Khi socket
đóng, nó thử lại sau 3 giây. Cleanup hủy timer, tháo handler và đóng socket, tránh
tạo nhiều kết nối khi người dùng đổi nguồn.

\subsection{Các ngưỡng thời gian}

\begin{table}[H]
\centering
\caption{Ngưỡng và giới hạn realtime}
\begin{tabularx}{\textwidth}{p{5cm}p{3cm}X}
\toprule
\textbf{Hằng số} & \textbf{Giá trị} & \textbf{Ý nghĩa}\\
\midrule
\codeword{DEVICE_STALE_MS} & 15 s & Không có device message thì root status chuyển waiting/offline.\\
\codeword{NODE_STALE_MS} & 20 s & Dùng thống kê online/offline.\\
\codeword{STALE_MS_NODES} & theo context & Dùng field online trong danh sách node.\\
\codeword{THROUGHPUT_WINDOW_SECONDS} & 60 s & Cửa sổ đồ thị bytes/kbps.\\
\codeword{MAX_DEBUG_LOGS} & 120 & Giới hạn log trong RAM.\\
Reconnect & 3 s & Thời gian mở lại frontend WS.\\
Server heartbeat & 30 s & Ping/pong và loại zombie socket.\\
\bottomrule
\end{tabularx}
\end{table}

Có nhiều hằng stale khác nhau giữa context và hook cũ
\codeword{useMeshNodesFromWebSocket} (7 giây). Hiện các hook ưu tiên trả dữ liệu
từ context nếu provider tồn tại, nhưng code fallback còn giữ logic trùng lặp.
Nên gom các timeout về một module cấu hình để tránh cùng node nhưng màn hình khác
nhau hiển thị trạng thái khác nhau.

\subsection{Tính throughput và latency}

Mỗi message text được đếm số byte UTF-8 theo giây. Cứ mỗi giây, provider lấy
bucket của giây trước, giữ 60 điểm và tính:
\[
  \text{kbps}=\frac{\text{bytes}\times 8}{1000}.
\]
Latency được tính từ thời gian browser nhận trừ RTC trong payload. Giá trị bị ép
tối thiểu 0. Nếu node dùng thời gian mặc định năm 2012 hoặc RTC chưa đồng bộ,
latency rất lớn và không phản ánh trễ mạng; UI cần đánh dấu time invalid thay vì
coi đó là network latency.

\section{Packet loss và server metrics}

Backend giữ \codeword{seqStats} theo định danh node. Mỗi 10 giây:
\[
  expected=maxSeq-minSeq+1,\qquad
  loss=\frac{expected-count}{expected}\times100\%.
\]
Kết quả được broadcast bằng \codeword{packetloss_update} và thử persist một
packet hệ thống. Cách tính này giả định seq tăng đều, không xử lý trùng gói và
reset/wrap sequence đầy đủ; ngưỡng \codeword{expected>10000} chỉ là bảo vệ sơ bộ.

Mỗi 5 giây, nếu có ít nhất một WebSocket client, backend dùng
\codeword{systeminformation} tạo \codeword{server_metrics}: CPU load, RAM,
nhiệt độ CPU, uptime và received time. Đây là metrics của \textbf{máy chạy
Node.js}, không phải MeshGateWay. Frontend giữ series riêng cho server và gateway.

\begin{figure}[H]
\centering
\resizebox{.92\textwidth}{!}{
\begin{tikzpicture}[node distance=9mm and 12mm]
 \node[block,minimum width=3.4cm] (seq) {Mesh packets\\seq theo node};
 \node[block,right=of seq,minimum width=3.5cm] (calc) {Cửa sổ 10 s\\tính loss};
 \node[block,right=of calc,minimum width=3.5cm,fill=green!8] (pl)
   {\codeword{packetloss_update}\\broadcast};
 \node[block,below=15mm of seq,minimum width=3.4cm] (host) {systeminformation\\máy backend};
 \node[block,right=of host,minimum width=3.5cm] (five) {Chu kỳ 5 s\\nếu có client};
 \node[block,right=of five,minimum width=3.5cm,fill=green!8] (sm)
   {\codeword{server_metrics}\\broadcast};
 \draw[dataflow](seq)--(calc);
 \draw[dataflow](calc)--(pl);
 \draw[dataflow](host)--(five);
 \draw[dataflow](five)--(sm);
\end{tikzpicture}}
\caption{Hai luồng số liệu do backend tự sinh}
\end{figure}

\section{SQLite và MongoDB: mức tương thích thực tế}

\begin{table}[H]
\centering
\caption{So sánh adapter hiện tại}
\begin{tabularx}{\textwidth}{p{4.3cm}p{4.4cm}X}
\toprule
\textbf{Khả năng} & \textbf{SQLite} & \textbf{MongoDB}\\
\midrule
Định danh \codeword{M} & Có, ưu tiên \codeword{M} & Không; validator và builder chỉ dùng \codeword{i}\\
Persist mesh & Có & Có nếu payload có \codeword{i}\\
Persist gateway status & Có & Không export hàm tương ứng\\
Node snapshot & \codeword{mac} & \codeword{ip}\\
History metadata & \codeword{getHistoryMacs} & \codeword{getHistoryIps}\\
History series input & \codeword{mac,sensorName,field} & \codeword{ip,field}\\
CSV export & Có & Không\\
Phù hợp server.js hiện tại & Có & Chưa tương thích đầy đủ\\
\bottomrule
\end{tabularx}
\end{table}

\textbf{Kết luận:} mặc dù \codeword{DB_TYPE} mặc định trong server.js là
\codeword{mongodb}, REST handler hiện gọi interface của SQLite. Chạy bằng
\codeword{execute.bat} an toàn hơn vì script đặt \codeword{DB_TYPE=sqlite}.
Muốn dùng MongoDB phải cập nhật adapter theo interface thống nhất hoặc tạo lớp
contract chung và test cho cả hai backend.

\section{Khởi chạy và dừng trên Windows}

\subsection{execute.bat}

\begin{figure}[H]
\centering
\resizebox{.86\textwidth}{!}{
\begin{tikzpicture}[node distance=7mm]
 \node[block,minimum width=8cm] (cd) {Chuyển tới thư mục chứa script};
 \node[decision,below=of cd] (pid) {File \codeword{.msms-dev.pid}\\trỏ tới process sống?};
 \node[block,right=22mm of pid,minimum width=4cm,fill=red!7] (stop)
   {Báo đang chạy\\không tạo instance mới};
 \node[block,below=15mm of pid,minimum width=8cm] (start)
   {Đặt \codeword{DB_TYPE=sqlite}\\Start-Process cmd /k npm run dev};
 \node[store,below=of start,minimum width=3.5cm] (pidfile) {Ghi PID cmd\\\codeword{.msms-dev.pid}};
 \node[block,below=of pidfile,minimum width=8cm,fill=green!8] (running)
   {React :3000 + Node.js :9090};
 \draw[flow](cd)--(pid);
 \draw[warnflow](pid)--node[above,font=\scriptsize]{có}(stop);
 \draw[flow](pid)--node[left,font=\scriptsize]{không}(start);
 \draw[flow](start)--(pidfile);
 \draw[dataflow](pidfile)--(running);
\end{tikzpicture}}
\caption{Cơ chế execute.bat}
\end{figure}

\subsection{stop.bat}

\codeword{stop.bat} yêu cầu quyền Administrator, đọc PID đã lưu và chạy
\codeword{taskkill /PID <pid> /T /F}. Sau đó script quét local endpoint cổng
3000 và 9090 bằng \codeword{netstat} để dọn process mồ côi. Việc dùng
\codeword{/T} là quan trọng vì PID file là PID của \codeword{cmd.exe}, còn React
và Node là tiến trình con.

\begin{lstlisting}[style=textstyle,caption={Cách sử dụng}]
cd /d "D:\Engineering Project\DashboardManagerWifiMesh"
execute.bat

rem Stop the program:
stop.bat
\end{lstlisting}

\section{Những phần giao diện chưa nối hoàn chỉnh}

Không phải mọi nút trong UI đều có backend tương ứng:
\begin{itemize}
 \item các nút Save changes, Reset to default của system setting và nhiều System
 Actions hiện chủ yếu là giao diện, chưa gọi API;
 \item thông tin System Uptime, Firmware/Hardware Version, Last Restart và Memory
 Usage trong card System Information có giá trị mẫu viết cứng;
 \item các trang OTA/FOTA có giao diện nhưng firmware gateway hiện chưa nối
 pipeline OTA hoàn chỉnh;
 \item notification email và các nút restart/factory reset chưa có endpoint
 điều khiển thực tế trong server.js.
\end{itemize}

Khi viết tài liệu vận hành cần phân biệt “màn hình tồn tại” với “hành động đã
được backend và firmware thực thi”.

\section{Rủi ro và đề xuất cải tiến}

\begin{enumerate}
 \item \textbf{Chuẩn hóa Store interface:} định nghĩa contract chung
 \codeword{connect}, \codeword{persistMeshPacket},
 \codeword{persistGatewayStatus}, \codeword{getNodeSnapshot},
 \codeword{getHistoryMacs}, \codeword{getHistorySeries} và export CSV.
 \item \textbf{Thống nhất định danh M:} đổi MongoDB từ field \codeword{ip} sang
 \codeword{mac} hoặc hỗ trợ migration/alias rõ ràng.
 \item \textbf{Giảm parser trùng lặp:} parser UART/mesh đang có ở server,
 context và hai hook. Tách shared package hoặc tạo bộ fixture contract.
 \item \textbf{Bảo vệ WebSocket và REST:} hiện chưa thấy authentication,
 authorization, rate limit hay giới hạn kích thước message ở lớp ứng dụng.
 \item \textbf{Giới hạn log:} server log toàn bộ text nhận được; telemetry dày
 có thể làm terminal chậm và lộ dữ liệu. Nên có log level và truncate.
 \item \textbf{Database unavailable:} server tiếp tục chạy sau lỗi connect,
 nhưng các lệnh store tiếp theo có thể liên tục lỗi. Nên expose health endpoint
 và trạng thái degraded.
 \item \textbf{CSV escaping:} dùng thư viện CSV hoặc escape đúng RFC.
 \item \textbf{RTC validity:} nhận biết năm mặc định/RTC chưa sync trước khi
 tính latency và lưu device time.
 \item \textbf{Packet loss:} xử lý duplicate, reset và wrap sequence bằng
 sliding window theo node.
 \item \textbf{CORS và network exposure:} backend bind 0.0.0.0 và CORS mở; cần
 cấu hình origin allowlist khi triển khai ngoài mạng lab.
\end{enumerate}

\section{Kế hoạch kiểm thử}

\begin{longtable}{p{4.3cm}p{5cm}p{4.7cm}}
\caption{Ma trận kiểm thử đề xuất}\\
\toprule
\textbf{Nhóm} & \textbf{Ca kiểm thử} & \textbf{Kết quả mong đợi}\\
\midrule\endfirsthead
\toprule
\textbf{Nhóm} & \textbf{Ca kiểm thử} & \textbf{Kết quả mong đợi}\\
\midrule\endhead
Parser & Mesh trực tiếp có \codeword{M} & Parse node, sensor và persist SQLite.\\
Parser & \codeword{uart_rx.payload} string & Parse đúng JSON lớp trong.\\
Parser & Hex bị chia nhiều frame & Carry ghép đúng một dòng.\\
Parser & Row có string/NaN/type lạ & Trả lỗi có kiểm soát, server không chết.\\
Realtime & Ngắt backend & UI chuyển wsOpen false, reconnect 3 s.\\
Realtime & Node im lặng 20 s & Thống kê offline cập nhật.\\
Config & Local/server/custom & Socket cũ đóng, socket mới mở; refresh vẫn giữ URL.\\
SQLite & Nhiều writer/read history & WAL cho phép truy vấn khi đang ghi.\\
History & from/to/limit & Thứ tự tăng dần và giới hạn đúng.\\
CSV & Hàng dữ liệu lớn & Stream theo chunk, file không thiếu dòng.\\
WebSocket & Client không pong & Server terminate sau heartbeat.\\
Windows & execute hai lần & Lần hai báo PID đang chạy.\\
Windows & stop sau PID mồ côi & Dọn process theo cổng 3000/9090.\\
MongoDB & Chạy toàn bộ API & Hiện dự kiến thất bại; dùng test để dẫn dắt sửa adapter.\\
\bottomrule
\end{longtable}

\section{Build và triển khai}

\begin{lstlisting}[style=textstyle,caption={Cài đặt và chạy development}]
cd DashboardManagerWifiMesh
npm install
npm install --prefix websocket-server
npm run dev
\end{lstlisting}

\begin{lstlisting}[style=textstyle,caption={Chạy riêng và build frontend}]
npm run start:react
npm run ws-server
npm run build
\end{lstlisting}

\begin{table}[H]
\centering
\caption{Biến môi trường quan trọng}
\begin{tabularx}{\textwidth}{p{4.7cm}p{4.4cm}X}
\toprule
\textbf{Biến} & \textbf{Mặc định} & \textbf{Ý nghĩa}\\
\midrule
\codeword{WS_PORT} / \codeword{PORT} & 9090 & Cổng HTTP + WebSocket.\\
\codeword{WS_HOST} & \codeword{0.0.0.0} & Interface bind.\\
\codeword{DB_TYPE} & mongodb trong server & Adapter database; execute.bat ghi đè sqlite.\\
\codeword{MONGO_URI} & local MeshData & URI MongoDB.\\
\codeword{SQLITE_DB_PATH} & mesh-data.sqlite & File SQLite.\\
\codeword{MDNS_HOSTNAME} & systemmsems & Hostname mDNS.\\
\codeword{REACT_APP_WS_URL} & local URL & URL frontend được đóng vào build.\\
\bottomrule
\end{tabularx}
\end{table}

\section{Kết luận}

Dự án đã có kiến trúc realtime hợp lý: một provider WebSocket dùng chung, backend
broadcast bản tin gốc, SQLite chuẩn hóa dữ liệu theo field và REST API phục vụ
history. Local mode hiện chạy cổng 9090, lưu SQLite và quảng bá
\codeword{systemmsems.local}. Những việc ưu tiên trước khi triển khai ổn định là
chuẩn hóa adapter MongoDB, loại bỏ parser trùng lặp, bảo vệ API/WS, phân biệt
chức năng UI mẫu với hành động thật và bổ sung test contract cho schema
\codeword{M}.

\appendix
\section{Thứ tự đọc mã đề xuất}

\begin{longtable}{p{6.3cm}p{7.7cm}}
\toprule
\textbf{Tệp} & \textbf{Nội dung nên đọc}\\
\midrule\endfirsthead
\toprule
\textbf{Tệp} & \textbf{Nội dung nên đọc}\\
\midrule\endhead
\pathcode{package.json} & Script chạy đồng thời frontend/backend.\\
\pathcode{src/index.js} & Cây provider.\\
\pathcode{src/context/meshRealtime.js} & Luồng realtime trung tâm.\\
\pathcode{src/utils/wsConfig.js} & Local/server/custom URL.\\
\pathcode{src/utils/meshUdpJsonSchema.js} & Schema và sensor mapping.\\
\pathcode{src/routes.js} & Route và menu.\\
\pathcode{src/layouts/system-settings/index.js} & Giao diện đổi nguồn dữ liệu.\\
\pathcode{websocket-server/server.js} & Toàn bộ HTTP/WS pipeline.\\
\pathcode{websocket-server/sqlite-store.js} & Database mặc định thực tế.\\
\pathcode{websocket-server/mongodb-store.js} & Adapter cũ cần đồng bộ.\\
\pathcode{execute.bat} & Khởi chạy và PID file.\\
\pathcode{stop.bat} & Dừng cây process và dọn cổng.\\
\bottomrule
\end{longtable}

\end{document}
